<<<P001629>%
<<<P001789><<<P001305>>
Подумайте о назначении <<<P000004> из формы <<<P001306>.
<<<P000008> Статический тип <<<P000009>.
Если <<<P000010> <<<P001790>, дальнейшие проверки не проводятся.
В противном случае это ошибка времени компиляции, если
<<<p000011> Метод получил название <<<P001307>.
<<<P000012> быть статичным типом
первый формальный параметр метода <<<P001308>,
<<<P000013> Статический тип второго.
Это ошибка времени компиляции, за исключением статического типа <<<P000014> <<<P000015>>
может быть присвоено <<<P000016> соответственно <<P000017>.
Является ли <<<P000018> <<<P001791>
Статический тип <<<P000019> Статический тип <<<P000020>.

<<<P001630>%
Оценка уступки <<<P000021> <<<P001309>
идет следующим образом:
Оценить <<<P000025> на объект <<<P000026>,
затем оценить <<P000027> для объекта <<<P000028>,
И, наконец, оценить <<P000029> для объекта <<<P000030>.
Назовите метод <<<P001310> на <<<P000031>
<<<P000032>> В качестве первого аргумента: <<P000033> как второй аргумент.
Далее <<<P000034> Показатель <<<P000035>.
% Стоит ли добавлять: Это динамическая ошибка, если <<P000036> оценивать на
Постоянный список или карта?
<<P001792>

<<<P001631>%
<<<P001793><<<P001311>>
Подумайте о назначении <<<P000039> формы <<<P001312>.
<<<P000042> Будьте суперклассом немедленного замкнутого класса.
Ошибка времени компиляции <<<P000043> не имеет метода <<<P001313>.
В противном случае пусть <<<P000044> быть статичным типом
первый формальный параметр метода <<<P001314>,
<<<P000045> Статический тип второго.
Это ошибка времени компиляции, если статический тип <<<P000046> соответственно <<P000047>
не может быть присвоено <<<P000048> соответственно <<P000049>.
Статический тип <<<P000050> Статический тип <<<P000051>.

<<<P001632>%
Для оценки присваивается форма <<<P001315>
Это эквивалентно выражению <<<P001316>.
<<<P001794>


<<<P001795> Назначение соединения
<<P001766>

<<<P001633>%
<<<P001796><<<P001317>>
Рассмотрим сложное назначение <<<P000058> <<<P001318>
где <<<P000061> является идентификатором или идентификатором, квалифицированным префиксом импорта.
Ошибки времени компиляции, которые могут быть вызваны
<<P001319>
Они также генерируются в случае <<<P000064>.
Статический тип <<<P000065> это
наименьшая верхняя граница статического типа <<<P000066> и статический тип <<<P000067>.

<<<P001634>%
Оценка комплексного назначения <<<P000068> <<<P001320>
идет следующим образом:
Оценить <<<P000071> на объект <<<P000072>.
Если <<P000073> не является нулевым объектом (<<<P001231>), <<P000074> Показатель <<<P000075>.
В противном случае оценка <<<P001321> к объекту <<<P000078>,
А потом <<<p000079> Оценка до <<<P000080>.
<<P001797>

<<<P001635>%
\Case{\code{<<<P000081>>>.<<<P000082>>> ??= <<<P000083>>>}>
Рассмотрим сложное назначение <<P000084> формы <<P001323>>>
где <<P000088> Тип буквальный
которые могут быть или не быть квалифицированы префиксом импорта.
Ошибки времени компиляции, которые могут быть вызваны
<<P001324>
Они также генерируются в случае <<<P000092>.
Статический тип <<<P000093> является наименьшей верхней границей
Статический тип <<<P001325> и статический тип <<<P000096>.

<<<P001636>%
Оценка комплексного назначения <<<P000097> формы \code{<<<P000098>>>.<<<P000099>>> ??= <<<P000100>>>}
где <<P000101> является типом буквальных поступлений следующим образом:
Оценка <<<P001327> Объект <<<P000104>.
Если <<P000105> не является нулевым объектом (<<<P001232>), <<P000106> Показатель <<<P000107>.
В противном случае оценка <<P001328> для объекта <<<P000111>,
<<<P000112>> Показатель <<<P000113>.
<<P001799>

<<<P001637>%
\Case\code{<<<P000114>>>.<<<P000115>>> ??= <<<P000116>>>}
Рассмотрим сложное назначение <<<P000117> из формы <<<P001330>.<<<P000121> Статический тип <<<P000122> <<<P000123>> Будьте свежей переменной типа <<P000124>.
За исключением ошибок внутри <<<P000125> и ссылки на имя <<<P000126>,
Ошибки времени компиляции, которые могут быть вызваны
<<P001331>
Они также генерируются в случае <<<P000130>.
Более того, это ошибка времени компиляции, если <<P000131> не имеет геттера с именем <<<P000132>.
Статический тип <<<P000133> является наименьшей верхней границей
Статический тип <<<P001332> и статический тип <<<P000136>.

<<<P001638>%
Оценка комплексного назначения <<P000137> формы <<P001333>>>
идет следующим образом:
Оценить <<<P000141> на объект <<<P000142>.
<<<P000143> быть свежей переменной, связанной с <<<P000144>.
Оценить <<<P001334> на объект <<<P000147>.
Если <<P000148> не является нулевым объектом (<<<P001233>), <<P000149> Показатель <<<P000150>.
В противном случае оценка <<<P001335> к объекту <<<P000154>,
<<<P000155>> Показатель <<<P000156>.
<<<P001801>

<<<P001639>%
\Case{<<<P001336>>
Рассмотрим сложное назначение <<<P000160> формы <<<P001337>.
Ошибки времени компиляции, которые могут быть вызваны
<<P001338>
Они также генерируются в случае <<<P000167>.
Кроме того, это ошибка времени компиляции
если статический тип <<<P000168> Не имеет значения <<<P001339>.
Статический тип <<<P000169> является наименьшей верхней границей
Статический тип <<<P001340> и статический тип <<<P000172>.

<<<P001640>%
Оценка комплексного назначения <<P000173> в форме
<<P001341>
идет следующим образом:
Оценить <<<P000177> на объект <<<P000178> а затем оценить <<P000179> Объект <<<P000180>.
Позвоните по телефону <<<P001342> Метод <<<P000181> с аргументом <<P000182>,
<<<P000183>> быть возвращенным объектом.
Если <<P000184> не является нулевым объектом (<<<P001234>), <<P000185> Показатель <<<P000186>.
В противном случае оценка <<P000187> для объекта <<<P000188>
Затем нажмите <<<P001343> Метод <<<P000189>
<<<P000190> В качестве первого аргумента: <<P000191> как второй аргумент.
<<<P000192> Показатель <<<P000193>.
<<<P001803>

<<<P001641>%
\Case{<<<P001344>>
Рассмотрим сложное назначение <<<P000196> из формы <<<P001345>.
Ошибки времени компиляции, которые могут быть вызваны
<<P001346>
Они также генерируются в случае <<<P000201>.
Кроме того, точно такие же ошибки времени компиляции, которые могут быть вызваны
Оценка выражения <<<P001347>
Они также генерируются в случае <<<P000203>.
Статический тип <<<P000204> является наименьшей верхней границей
Статический тип <<<P001348> и статический тип <<<P000206>.

<<<P001642>%
Оценка комплексного назначения <<<P000207> <<<P001349>>
идет следующим образом:
Оценить <<<P001350> на объект <<<P000211>.
Если <<P000212> не является нулевым объектом (<<<P001235>), затем <<<P000213> Показатель <<<P000214>.
В противном случае оценка <<<P001351> <<<P000217>,
<<<P000218>> Показатель <<<P000219>.
<<<P001805>

<<<P001643>%
\Case{<<<P001352>>
Рассмотрим сложное назначение <<P000223> формы <<<P001353>.
Ошибки времени компиляции, которые могут быть вызваны
<<P001354>
Они также генерируются в случае <<<P000230>.
% Примечание: Мы используем статический тип <<<P001355> вместо того, чтобы
% <<<P001356> Хотя последнее было бы проще. Это потому, что
% первое останется верным, если будет введено NNBD, и поскольку оно уменьшает
% от количества синтетического синтаксиса.
Статический тип <<<P000235> является наименьшей верхней границей
Статический тип <<<P001357> и статический тип <<<P000238>.

<<<P001644>%
Оценка комплексного назначения <<<P000239> формы \code{<<<P000240>>>?.<<<P000241>>> ??= <<<P000242>>>}
идет следующим образом:
Оценить <<<P000243> на объект <<<P000244>.Если <<P000245> является нулевым объектом (<<<P001236>), затем <<<P000246> Оценивает нулевой объект.
В противном случае пусть <<<P000247> быть свежей переменной, связанной с <<<P000248>.
Оценка <<<P001359> Объект <<<P000251>.
Если <<P000252> не является нулевым объектом (<<<P001237>), затем <<<P000253> Показатель <<<P000254>.
В противном случае оценка <<<P001360> к объекту <<<P000258>,
<<<P000259>> Показатель <<<P000260>.
<<<p001807>

<<<P001645>%
<<<P001808><<<P001361>>
Составное назначение формы <<<P001362>
где <<P000267> Тип буквальный
который может или не может быть квалифицирован импортным префиксом
Это эквивалентно выражению <<<P001363>.
<<<P001809>

<<<P001646>%
\Case{<<<P001364>>
Для любого другого действующего оператора <<P000274>,
a сложное присвоение формы <<P001365>
эквивалентен <<<P001366>,
где <<P000282> является идентификатором или идентификатором, квалифицированным префиксом импорта.
<<<P001811>

<<<P001647>%
\Case{<<P001367>>
Составное назначение формы <<<P001368>
где <<P000291> Тип буквальный
который может или не может быть квалифицирован импортным префиксом
Это эквивалентно P001369.
<<<P001813>

<<<P001648>%
\Case{<<<P001370>>
Рассмотрим сложное назначение <<P000302> формы <<<P001371>.
<<<P000307> Свежие переменные, статическим типом которых является статический тип <<<P000308>.
За исключением ошибок внутри <<<P000309> и ссылки на имя <<<P000310>,
Ошибки времени компиляции, которые могут быть вызваны
<<P001372>
Они также генерируются в случае <<<P000317>.
Статический тип <<<P000318> Статический тип <<<P001373>.

<<<P001649>%
Оценка комплексного назначения <<P000323> формы <<P001374>
идет следующим образом:
Оценить <<<P000328> на объект <<<P000329> <<<P000330>> быть свежей переменной, связанной с <<P000331>.
Оценить <<<P001375> на объект $r$
<<<P000339>> Показатель <<<P000340>.
<<<P001815>

<<<P001650>%
\Case{<<<P001376>>
Рассмотрим сложное назначение <<P000345> формы <<<P001377>.
Пусть <<<P000350> и <<<P000351> Будьте свежими переменными
где статическим типом первого является статический тип <<<P000352>
статическим типом последнего является статический тип <<<P000353>.
За исключением ошибок внутри <<P000354> и <<P000355> и ссылки на
имена <<P000356> и <<P000357>,
Ошибки времени компиляции, которые могут быть вызваны
<<P001378>
Они также генерируются в случае <<<P000364>.
Статический тип <<<P000365> Статический тип <<<P001379>.

<<<P001651>%
Оценка комплексного назначения <<<P000370>> Форма
<<P001380>
идет следующим образом:
Оценить <<<P000375> на объект <<<P000376> и оценить <<P000377> Объект <<P000378>.
<<<P000379> и <<<P000380> быть свежими переменными, связанными с <<<P000381> и <<<P000382> соответственно.
Оценить <<<P001381> на объект <<<P000389>,
А потом <<P000390> Показатель <<P000391>.
<<<P001817>

<<<P001652>%
\Case\code{<<<P000392>>>?.<<<P000393>>> <<<P000394>>>= <<<P000395>>>}
Рассмотрим сложное назначение <<P000396> формы <<<P001383>.
Ошибки времени компиляции, которые могут быть вызваны
<<P001384>
Они также генерируются в случае <<<P000405>.
Статический тип <<<P000406> Статический тип <<<P001385>.

<<<P001653>%
Оценка комплексного назначения <<<P000411> в форме
<<P001386>
идет следующим образом:
Оценить <<<P000416> на объект <<<P000417>.
Если <<P000418> является нулевым объектом, то <<<P000419> оценивает нулевой объект (<<<P001238>).
В противном случае пусть <<<P000420> быть свежей переменной, связанной с <<<P000421>.
Оценить <<<P001387> на объект <<<P000426>.
Далее <<<P000427> Показатель <<<P000428>.
<<<P001819>

<<<P001654>%
<<<P001820><<<P001388>>
Составное назначение формы <<<P001389>
где <<P000437> Тип буквальныйЭто эквивалентно выражению <<<P001390>.
<<<P001821>


<<<<P001822>{условный}
<<P001767>

<<<P001655>%
<<<P001286> Оценка одного из двух выражений
основанный на боолеиновом состоянии.

<<P001173>
<условный Выражение>::= <ifNullExpression>
\gnewline{} (?<expressionWithoutCascade>: <expressionWithoutCascade>)?
<<<p001202>

<<<P001656>%
Оценка условного выражения <<P000442> формы <<<P000443>
идет следующим образом:

<<<P001657>%
Во-первых, <<P000444> оценивается на предмет <<P000445>.
% Эта ошибка может произойти из-за неявных отливок и нулевых.
Это динамическая ошибка, если тип времени выполнения <<<P000446> Это не P001391.
Если <<<P000447> <<<P001824>, то значение <<<P000448> это
Результат оценки выражения <<<P000449>.
В противном случае значение <<<P000450> является результатом оценки выражения <<<P000451>.

<<<P001658>%
Если <<P000452> показывает локальную переменную <<<P000453> имеет тип <<<P000454>,
Тип <<<P000455> Известно, что <<<P000456> в <<<P000457>,
Если ни одно из следующих утверждений не верно:

<<P001174>
<<<P001825> <<P000458> потенциально мутирует в <<<P000459>,
<<<P001826> <<P000460> потенциально мутирует в пределах функции
, где <<<P000461> объявляется, или
<<<P001827> <<P000462> доступ к функции, определенной в <<<P000463> и
<<<P000464> потенциально мутирует в любом месте в области <<<P000465>.
<<<P001203>

<<<P001659>%
Ошибка времени компиляции, если
Статический тип <<<P000466> Не может быть присвоено <<<P001392>.
Статический тип <<<P000467> наименьшая верхняя граница (<<<P001239>) из
Статический тип <<<P000468> и статический тип <<<P000469>.


<<<P001828> Если нулевые выражения
<<P001768>

<<<P001660>%
<<<P001287> оценивает выражение и
Если результатом является нулевой объект (<<<P001240>), оценивает другой.

<<P001175>
<ifNullExpression> ::= <logicalOrExpression> (??? <logicalOrExpression>)
<<<P001204>

<<<P001661>%
Оценка выражения <<P000470> формы <<P001393>
идет следующим образом:

<<<P001662>%
Оценить <<<P000473> на объект <<<P000474>.
Если <<P000475> не является нулевым объектом (<<<P001241>), затем <<<P000476> Показатель <<<P000477>.
В противном случае оценка <<<P000478> к объекту <<<P000479>,
А потом <<P000480> Показатель <<P000481>.

<<<P001663>%
Статический тип <<<P000482> наименьшая верхняя граница (<<<P001242>) из
Статический тип <<<P000483> и статический тип <<<P000484>.


<<<P001829> Логические булевы выражения
<<P001769>

<<<P001664>%
Логические булевы выражения объединяют булевы объекты с помощью
Булевы операторы соединения и разъединения.

<<P001176>
<logicalOrExpression> ::= \gnewline{}
<logicalAndExpression> (<logicalAndExpression>)

<логический AndExpression> ::= <equalityExpression> (<<<P002346><<<P002347>> <equalityExpression>)
<<P001205>

<<<P001665>%
<<<P001288> является выражением равенства
(<<<P001243>),
или вызов логического булевого оператора на
Выражение <<P000485> с аргументом <<P000486>.

<<<P001666>%
Оценка логического булевого выражения <<P000487> формы <<<P000488> причины
Оценка <<<P000489> Объект <<<P000490>.
% Эта ошибка может произойти из-за неявных отливок и нулевых.
Это динамическая ошибка, если тип времени выполнения <<<P000491> Это не P001394.
Если <<P000492> <<<P001831>, то результат оценки <<P000493> <<<P001832>,
иначе <<P000494> оценивается на предмет <<P000495>.
% Эта ошибка может произойти из-за неявных отливок и нулевых.
Это динамическая ошибка, если тип времени выполнения <<<P000496> Это не P001395.
В противном случае результат оценки <<<P000497> является <<<P000498>.

<<<P001667>%
Оценка логического булевого выражения <<<P000499> формы <<<P000500>Причин оценки <<<P000501> Производить объект <<P000502>.
% Эта ошибка может произойти из-за неявных отливок и нулевых.
Это динамическая ошибка, если тип времени выполнения <<<P000503> Это не P001396.
Если <<<P000504> <<<P001833>, то результат оценки <<P000505> <<<P001834>,
иначе <<P000506> Оценивается на объект <<<P000507>.
% Эта ошибка может произойти из-за неявных отливок и нулевых.
Это динамическая ошибка, если тип времени выполнения <<<P000508> Это не P001397.
В противном случае результат оценки <<<P000509> является <<<P000510>.

<<<P001668>%
Логическое булево выражение <<P000511> формы <<P000512>>>
показывает локальную переменную <<<P000513> имеет тип <<<P000514>
если выполняются оба следующих условия:

<<P001177>
<<<P001835> Либо <<<P000515> Показывает, что <<<P000516> имеет тип <<<P000517>
<<<P000518> Показывает, что <<<P000519> имеет тип <<P000520>.
<<<P001836> <<P000521> не мутирует в <<<P000522> или в пределах функции
За исключением того, где объявлен <<P000523>.
<<<P001206>

<<<P001669>%
Если <<P000524> показывает локальную переменную <<<P000525> имеет тип <<P000526>,
Тип <<<P000527> Известен как <<<P000528> в <<<P000529>,
Если ни одно из следующих утверждений не верно:

<<P001178>
% % Первый пункт здесь не нужен для звука,
Сохраняется в основном для обратной совместимости.
% % Если <<P000530> Показывает, что <<<P000531> имеет тип <<<P000532>, затем любое назначение
% % в <<<P000533> Это не является недействительным.
Удаление линии видно в семантике, потому что:
% % от числа x;
%% (x??=42) null & x is int & x.toRadixString(16)!=""
% допускается без линии и не допускается.
% % На момент написания анализатор запрещает код.
<<<P001837> <<P000534> потенциально мутирует в <<<P000535>,
<<<P001838> <<P000536> потенциально мутирует в <<<P000537>,
<<<P001839> <<P000538> потенциально мутирует в пределах функции
где <<<P000539> объявляется, или
<<<P001840> <<P000540> доступ к функции, определенной в <<<P000541> и
<<P000542> Потенциально мутирует в любом месте в области <<<P000543>.
<<P001207>

<<<P001670>%
Ошибка времени компиляции, если
Статический тип <<<P000544> не может быть присвоено <<<P001398>
или если статический тип <<<P000545> не может быть назначено на <<<P001399>.
Статический тип логического булевого выражения <<<P001400>.


<<<P001841>{Равенство}
<<<P001770>

<<<P001671>%
Выражения равенства проверяют объекты на равенство.

<<P001179>
<equalityExpress> ::= <<<P001842><
<relationalExpression> (<equalityOperator> <relationalExpression>)
<<P001843> \SUPER{}<equalityOperator> <Реляционное выражение>

<equalityOperator> ::=== "
<<<P001845 >>!= "
<<<P001208>

<<<P001672>%
<<<P001289> является либо реляционным выражением
(<<P001244>),
или вызов оператора равенства на любой <<<P001846>
или выражение <<<P000546>, с аргументом <<P000547>.

<<<P001673>%
Оценка выражения равенства <<P000548> формы <<P001401>>>
идет следующим образом:

<<P001180>
<<P001847> Выражение <<P000551> Оценивается на объект <<<P000552>.
<<<P001848> Выражение <<P000553> оценивается на предмет <<P000554>.
<<<P001849> Если <<<P000555> или <<P000556> является нулевым объектом (<<<P001245>),
<<<P000557> Показатель <<<P001850>{} <<<P000558> и <<<P000559> является нулевым объектом
и <<<P001851>{} в противном случае.
В противном случае
<<P001852> Оценка <<<P000560> эквивалентно методу вызова
<<<P001402>.
<<<p001209>

<<<P001674>%
Оценка выражения равенства <<P000563> в форме
\code{<<P001854>>>{== <<P000564>>
идет следующим образом:

<<P001181>
<<<P001855> Выражение <<P000565> оценивается на предмет <<P000566>.<<<P001856> Если \THIS{{} или <<P000567> является нулевым объектом (<<<P001246>),
<<<P000568> Оценки для оценок \TRUE{}
Если оба \THIS{{} и $o$ является нулевым объектом
и <<<P001860>{} в противном случае.
В противном случае
<<<P001861> Оценка <<<P000570> эквивалентно методу вызова
<<<P001403>.
<<<p001210>

<<<P001862>{%>
В результате приведенного выше определения,
Пользовательский <<<P001863>{=== Методы могут предполагать, что их аргумент не является нулевым.
Избегайте стандартного прелюдия котельной:

<<<P001404>

Еще один вывод заключается в том, что необходимости никогда не существует.
Использование <<<P001405> Для проверки на соответствие <<<P001864>,
Никто не должен беспокоиться о том, писать ли
\NULL{{=<<P000572> или <<P000573>== <<<P001866>>%
?

<<<P001675>%
Выражение равенства формы <<<P001406> является эквивалентным
Выражение <<<P001407>.
Выражение равенства формы \code{\SUPER{}!= <<P000578> является эквивалентным
выражение <<<P001869><! (\SUPER{}) <<<P000579>>.

% выражение <<<P000580> оценивается на объект <<<P000581>;
% от выражения <<P000582> оценивается на предмет <<P000583>.
% Далее, если $o_1$ и $o_2$ Это один и тот же объект,
% далее <<<P000586> <<<P001871>, в противном случае <<P000587>>> Показатель <<<P001872>.

<<<P001676>%
Статический тип выражения равенства <<<P001408>.


<<<P001873>{Реляционные выражения}
<<<P001771>

<<<P001677>%
Реляционные выражения вызывают реляционных операторов на объектах.

<<P001182>
<реляционный Выражение>::= <bitwiseOrExpression> <<<<P001874><
(<typeTest> | <typeCast> | <relationalOperator> <bitwiseOrExpression>)?
<<<P001875> \SUPER{} <relationalOperator> <bitwiseOrExpression>

<relationalOperator> ::=>= "
<<<P001877>> "
<<<P001878 ><=
<<<P001879>< "
<<<P001211>

<<<P001678>%
<<<P001290> Это либо немногословное выражение
(<<<P001247>),
или вызов реляционного оператора на любой <<<P001880>
или выражение <<<P000588>, с аргументом <<P000589>.

<<<P001679>%
Реляционное выражение формы <<<P000590> <<P000591> <<P000592> является эквивалентным
Метод вызова <<<P001409>.
Реляционное выражение формы <<<P001881> <<P000596> <<P000597> является эквивалентным
Метод вызова <<<P001410>.


<<<<P001882>{Битвы выражений}
<<P001772>

<<<P001680>%
Битовые выражения вызывают битовые операторы на объектах.

<<P001183>
<bitwiseOrExpression> ::= <<<<P001883><
<bitwiseXorExpression> (<bitwiseXorExpression>)*
<<<P001884> \SUPER{} (' |' <bitwiseXorExpression>) +

<bitwiseXorExpress> ::= <<<<P001886><
<bitwiseAndExpression> ('^' <bitwiseAndExpression>)
\alt \SUPER{} ('^' <bitwiseAndExpression>)+

<bitwiseAndExpression> ::= <shiftExpression> (<<<P002348>> <shiftExpression>) *
<<<P001889> <<<P001890>< (<<<P002349>> <shiftExpression>)+

<bitwiseOperator> ::=<<<P002350> "
<<<P001891>> "
<<<P001892>> "
<<<p001212>

<<<P001681>%
<<<P001291> является выражением сдвига (<<<P001248>),
или обращение к оператору bitwise
В-третьих, на любой В.М.: В.М.: В.М.: В.М. В.М.: В.М., в частности, в связи с тем, что
с аргументом <<P000601>.

<<<P001682>%
Пошаговое выражение формы <<<P001411> является эквивалентным
Метод вызова <<P000605>.
Побитовое выражение формы \code{\SUPER{} <<P000606> <<P000607>> является эквивалентным
Метод вызова <<P001412>.

<<<P001896>{%>
Должно быть очевидно, что правила статического типа для этих выражений
Определяется эквивалентностью выше --ergo,
Правила типа для вызова метода и
подписи операторов по типу <<<P000610>.
То же самое относится и к аналогичным ситуациям в рамках данной спецификации.%
?


<<<P001897> Сдвиг!
<<p001773>

<<<P001683>%Выражения сдвига вызывают операторов сдвига на объектах.

<<P001184>
<shiftExpression> ::= \gnewline{}
<additiveExpression> (<shiftOperator> <additiveExpression>)*
<<<P001899> \SUPER{} (<shiftOperator> <additiveExpression>)+

<shiftOperator> ::=<<<P001901> "
<<<P001902><<<P001903> "
<<<P001904><<<P001905> "
<<<P001213>

<<<P001684>%
<<<P001292> является либо аддитивным выражением
(<<<P001249>),
или вызов оператора смены
либо <<<<P001906>< или выражение <<P000611>,
с аргументом <<P000612>.

<<<P001685>%
Сдвиг выражения формы <<<P000613> <<P000614> <<P000615> является эквивалентным
Метод вызова <<P001413>.
Сдвиг выражения формы <<<P001907> <<P000619> <<P000620> является эквивалентным
Метод вызова <<<P001414>.

<<<P001908>{%>
Обратите внимание, что это определение подразумевает порядок оценки слева направо
Среди сменных выражений:
<<<P001415>
оценивается как
<<<P001416> который эквивалентен
<<P001417>.
То же самое относится к аддитивным и мультипликативным выражениям.
?


<<<p001909>> Аддитивные выражения
<<P001774>

<<<P001686>%
Аддитивные выражения вызывают операторов сложения на объектах.

<<<<<<<<<<<<<<<<<<<<<<<<<<<<<<<<<<<<<<<<<<<<<<<<<<<<<<<<<<<<<<<<<<<<<<<<<<<<<<<<<<<<<<<<<<<<<<<<<<<<<<<<<<<<<<<<<<<<<<<<<<<<<<<<<<<<<<<<<<<<<<<<<<<<<<<<<<<<<<<<<<<<<<<<<<<<<<<<<<<<<<<<<<<<<<<<<<<<<<<<<<<<<<<<<<<<<<<<<<<<<<<<<<<<<<<<<<<<<<<<<<<<<<<<<<<<<<<<
<additiveExpression> ::= <multiplicativeExpression>
\gnewline{} (<additiveOperator> <multiplicativeExpression>)*
<<<p001911> \SUPER{} (<additiveOperator> <multiplicativeExpression>)+

<additiveOperator> ::=> "
<<<P001913>> "
<<<P001214>

<<<P001687>%
<<<P001293> является мультипликативным выражением
(<<<P001250>),
или вызов аддитивного оператора
либо <<<<P001914>< или выражение <<<P000632>,
с аргументом <<P000633>.

<<<P001688>%
Добавочное выражение формы <<<P000634> <<P000635><<P000636> является эквивалентным
Метод вызова <<<P001418>.
Добавочное выражение формы \SUPER{} <<P000640> <<P000641> является эквивалентным
Метод вызова <<<P001419>.

<<<P001689>%
Статический тип аддитивной экспрессии обычно определяется
по подписи, указанной в декларации используемого оператора.
Призывы операторов <<<P001420> и <<<P001421> из
класс <<<P001422>, <<<P001423> и <<<P001424>
Обрабатываются специально чекером.

<<<P001690>%
<<<P000644> быть аддитивным выражением формы <<<P001425>,
<<<P000648> Статический тип <<<P000649>,
<<<P000650>> Контекстный тип <<<P000651>.

Если \SubtypeNE{T}{<<P001426>> и не \SubtypeNE{T}{<<P001427>>>}
Тип контекста <<<P000652> определяется следующим образом:

<<P001186>
<<<<P001918>< Если \SubtypeNE{<<P001428>>>}{C}, а не <<P001920>>>{<<P001429>>>}{C},
и \SubtypeNE{T}{<<P001430>>>}}
Контекстный тип <<<P000653> является <<<P001431>.
<<<<P001922>< Если \SubtypeNE{<<P001432>>>}{C}, а не <<P001924>>>{<<P001433>>>}{C},
и не \SubtypeNE{T}{<<P001434>>>
Тогда контекстный тип <<<P000654> является <<<P001435>.
<<<<P001926>< В противном случае тип контекста <<<P000655> является <<<P001436>.
<<<P001215>

<<<P001691>%
Далее <<<P000656> Статический тип <<<P000657>.
Если <<<P001927>{T}{<<P001437>> и не \SubtypeNE{T}{\code{Never}>
<<<P000658> присваивается <<<P001439>,
Статический тип <<<P000659> определяется следующим образом:

<<P001187>
<<<P001929>< Если \SubtypeNE{T}{<<P001440>>
Статический тип <<<P000660> является <<<P000661>.
<<<<P001931>< В противном случае, если \SubtypeNE{S}{<<P001441>>>}
и не <<<P001933 >>{S}{<<<P001442>>,
Статический тип <<<P000662> является <<<P001443>.
<<<P001934>< В противном случае, если \SubtypeNE{T}{\code{int}},
\SubtypeNE{S}{<<P001445>>>}} и не \SubtypeNE{S}{<<P001446>>>,
Статический тип <<<P000663> является <<<P001447>.<<<<P001938>< В противном случае статический тип <<<P000664> является <<<P001448>.
<<<P001216>

<<<P001939> Множественные выражения
<<P001775>

<<<P001692>%
Умножительные выражения вызывают операторы умножения на объекты.

<<P001188>
<multiplicativeExpress> ::= <<<<P001940>>
<unaryExpression> (<multiplicativeOperator> <unaryExpression>)*
<<<p001941> \SUPER{} (<multiplicativeOperator> <unaryExpression>)+

<multiplicativeOperator> ::= '* "
<<<P001943>> "
<<<P001944><<<P002351> "
<<<P001945></> "
<<P001217>

<<<P001693>%
<<<P001294> является унарным выражением
(<<<P001251>),
или вызов мультипликативного оператора
либо <<<<P001946>< или выражение <<P000665>,
с аргументом <<P000666>.

<<<P001694>%
Умножительное выражение формы <<<P000667> <<P000668> <<<P000669> эквивалентно
Метод вызова <<P001449>.
Множественное выражение формы <<<P001947> <<P000673> <<P000674> является эквивалентным
Метод вызова <<P001450>.

<<<P001695>%
Статический тип мультипликативного выражения обычно определяется
по подписи, указанной в декларации используемого оператора.
Призывы операторов <<<P001451> и <<<P001452> из
класс <<<P001453>, <<<P001454> и <<<P001455>
Обрабатываются специально чекером.

<<<P001696>%
<<<P000677> быть мультипликативным выражением формы <<<P001456>
где <<P000681> является одним из <<<P001457> или <<<P001458>,
<<<P000682> Статический тип <<<P000683>,
<<<P000684>> Контекстный тип <<<P000685>.

Если <<<P001948>{T}{<<P001459>> и не \SubtypeNE{T}{<<P001460>>>}
Тип контекста <<<P000686> определяется следующим образом:

<<P001189>
<<<<P001950>< Если \SubtypeNE{<<P001461>>>}{C}, а не <<P001952>>>{<<P001462>>>}{C},
и \SubtypeNE{T}{<<P001463>>>}}
Контекстный тип <<<P000687> является <<<P001464>.
<<<<P001954>< Если \SubtypeNE{<<P001465>>>}{C}, а не <<P001956>>>{<<P001466>>>}{C},
и не <<<P001957>{T}{<<<P001467>>
Контекстный тип <<<P000688> является <<<P001468>.
<<<<P001958>< В противном случае контекстный тип <<<P000689> является <<<P001469>.
<<<P001218>

<<<P001697>%
Пусть далее <<<P000690> будет статичным типом <<<P000691>.
Если <<<P001959>{T}{<<P001470>> и не \SubtypeNE{T}{\code{Never}>
<<P000692> присваивается <<<P001472>,
Статический тип <<<P000693> определяется следующим образом:

<<P001190>
<<<<P001961>< Если <<<P001962>{T}{<<P001473>>
Статический тип <<<P000694> является <<<P000695>.
<<<P001963>< В противном случае, если <<P001964>>>{S}{<<P001474>>>}
и не <<<P001965 >>{S}{<<<P001475>>,
Статический тип <<<P000696> является <<<P001476>.
<<<P001966>< В противном случае, если \SubtypeNE{T}{\code{int}},
\SubtypeNE{S}{<<P001478>>>}} и не \SubtypeNE{S}{\code{Never},
Статический тип <<<P000697> является <<<P001480>.
<<<<P001970>< В противном случае статический тип <<<P000698> является <<<P001481>.
<<<P001219>


<<<P001971> Унарные выражения
<<<P001776>

<<<P001698>%
Унарные выражения вызывают унарных операторов на объектах.

<<P001191>
<unaryExpression> ::= <prefixOperator> <unaryExpression>
<<<P001972> <awaitExpression>
<<<P001973> <postfixExpression>
<<<P001974> (<minusOperator> | <tildeOperator>) <<<P001975><
<<<P001976> <incrementOperator> <assignableExpression>

<prefixOperator> ::= <minusOperator>
<<<<P001977> <отрицатель
<<<p001978> <tildeOperator>

<minusOperator> ::= "

<negationOperator> ::= '! "

<tildeOperator> ::=~ "
<<<P001220>

<<<P001699>%
<<<P001295> Это либо постфиксное выражение
(<<<P001252>),
Выражение ожидания (<<<P001253>)
или вызов оператора префикса на выражениеили вызов унарного оператора на любой <<<P001979> Выражение <<P000699>.

<<<p001700>%
Выражение <<P000700> рассматривается как
(<<<P001254>)
<<P001482>.

<<<P001701>%
Выражение формы <<<P001483> рассматривается как
<<P001484>.
Выражение формы \code{-{}-<<P000704>>>} трактуется как
<<P001485>.

<<<P001702>%
<<<P000706> быть выражением формы <<<P001486>
где <<P000708> является буквальным целым числом (<<P001255>) с числовым целым числом <<P000709>,
со статическим типом контекста <<P000710>.
Если <<<P001487> Применяется к <<<P000711> и <<<P001488> не присваивается <<<P000712>,
Статический тип <<P000713> является <<<P001489>;
Статический тип <<<P000714> является <<<P001490>.

<<<P001703>%
Если статический тип <<P000715> <<<P001491>, то <<<P000716> оценивать на
Пример <<<P001492> класс, представляющий числовое значение <<<P000717>.
Если <<P000718> Это ноль и P001493. класс может представлять собой отрицательное нулевое значение,
Вместо этого полученный экземпляр представляет собой отрицательное нулевое значение.
Ошибка времени компиляции, если целое число <<P000719> не могут быть представлены
Именно на примере <<<P001494>.

<<<P001704>%
Если статический тип <<P000720> <<<P001495>, то <<<P000721> оценивать на
Для примера <<<P001496> класс, представляющий числовое значение <<P000722>.
Если <<P000723> равна нулю, полученный экземпляр представляет собой
<<<P001981>{отрицательное} нулевое двойное значение, <<P001497>.
Это ошибка времени компиляции, если целое число <<<P000724> не могут быть представлены
Именно на примере <<<P001498>.
<<<P001982>{%>
Мы рассматриваем <<<P001499> <<<P001983>{как будто} это одно целое число
с отрицательным числовым значением.
Мы не оцениваем <<P000726> индивидуально, как выражение,
или заботится о своем статическом типе.
?

<<<P001705>%
Любое другое выражение формы <<<P001500> является эквивалентным
Метод вызова <<<P001501>.
Выражение формы <<<P001502> является эквивалентным
метод вызова (<<<P001256>) <<<P001503>.


<<<P001984> Подождите выражения.
<<P001777>

<<<P001706>%
<<<P001296> позволяет кодировать
Управление выходом до асинхронной операции
(<<<P001257>)
Готово.

<<P001192>
<awaitExpress> ::= <<<P001985><unaryExpression>
<<P001221>

<<<P001707>%
<<<P001786>%
<<<P000732> Выражение формы <<<P001504>.
<<<P000734> Статический тип <<<P000735>.
Статический тип <<<P000736> Затем <<<<P001986> С.
(<<P001258>).

<<<P001708>%
Оценка <<<P000737> идет следующим образом:
Во-первых, выражение <<P000738> оценивается на предмет <<P000739>.
Пусть \DefineSymbol{T} будет \flatten С.
Если тип времени выполнения <<<P000740> является подтипом <<<P001505>,
Тогда пусть \DefineSymbol{f} be $o$
В противном случае пусть <<<P000743> быть результатом создания
новый объект с использованием конструктора <<<P001506>
<<<P000745> В качестве аргумента.

<<<P001709>%
Далее, поток, связанный с
Внутренняя замкнутая асинхронная <<<P001990>< петля
(<<<P001259>),
Если есть, то приостанавливается.
Нынешнее призвание функционального тела, непосредственно заключающего в себе <<P000746>
Приостанавливается до тех пор, пока <<P000747> Готово.
Спустя некоторое время после завершения операции <<P000748> управление возвращается к текущему вызову.
Если <<P000749> Завершилась ошибка <<<P000750> и след стека <<<P000751>,
<<P000752> броски <<<P000753> и <<<P000754>
(<<<P001260>).
Если <<P000755> завершается объектом <<<P000756>, <<P000757> Показатель <<P000758>.

<<<<P001991>{%>
Использование \flattenName{} для поиска <<<P000759>
и, следовательно, определить динамический тип испытания
Это означает, что мы ожидаем будущего в каждом случае, когда этот выбор будет правильным.
?

<<<<P001993>{%>
Интересный случай на краю этого компромисса - когда <<P000760>Имеет статический тип <<<P001507>.
Можно сказать, что за этим типом стоит намерение
Значение <<<P000761> является <<<P001508>,
Или это <<<P001509> Который не является будущим,
Или это <<<P001994>.
Так что, предположительно, мы должны ждать первого рода,
Второе и третье должны быть без изменений.
Однако второй вид может быть <<<P001510>.
Этот объект не является <<<P001511>, и это не <<<P001995>,
Таким образом, он <<<P001996>{должен} считаться во второй группе.
Тем не менее, <<<P001997>\code{FutureOr<Object>?}><<<P001513>>
Таким образом, мы ожидаем <<<P001998> <<<P001514>.
Мы выбрали эту семантику, потому что это было наименьшее изменение.
относительно семантики в более ранних версиях Дарта,
А также потому, что позволяет простое правило:
Тип <<<P001515> Используется для того, чтобы решить, есть ли
Будущее (если оно есть) ждет, и исключений нет.
В таких случаях, как этот пример, тип подразумевает, что
<<<P001516> Не следует ждать.
Таким образом, мы ожидаем каждое будущее, которое мы можем разумно ожидать.
?

<<<<P001999>{%>
Выражение ожидания может происходить только в объявленной функции.
асинхронно. <<<P002000>> идентификатор не имеет особого значения
в контексте нормальной функции, поэтому случаи <<<P002001>
В этих функциях не вводится выражение ожидания.
Тем не менее, <<P001517> может быть действительным вызовом функции в
неасинхронные функции.%
?

<<<<P002002><%
Ожидаемое выражение не могло осмысленно возникнуть в синхронной функции.
Если бы такая функция была приостановлена в ожидании будущего,
Он больше не будет синхронным.
?

<<<P002003><%
Это не ошибка времени компиляции, если тип <<P000764> не является
супертип или подтип <<<P001518>.
Инструменты могут дать подсказку в таких случаях.%
?


<<<P002004> Постфиксные выражения
<<P001778>

<<<P001710>%
Выражения Postfix вызывают операторов postfix на объектах.

<<P001193>
<postfixExpression> ::= <assignableExpression> <postfixOperator>
<<P002005> <primary> <selector> *

<postfixOperator> ::= <IncrementOperator>

<constructorInvocation> ::= <<<P002006><
<typeName> <typeArguments> <identifier> <аргументы>

<selector> ::= '! "
<<P002007> <assignableSelector>
<<<P002008> <Аргументы>

<argumentpart> ::=
<тип Аргументы>. <аргументы>

<incrementOperator> ::=++ "
<<<P002009><-<<P002010>> "
<<<P001222>

<<<P001711>%
<<<P001297> является первичным выражением;
функция, метод или вызов Геттера;
Призыв названного конструктора;
или вызов оператора постфикса на выражение <<<P000765>.
Все, кроме последних двух, указаны в другом месте.

<<<P001712>%
<<<P002011> Призывы конструктора
Рассмотрим <<<<P002012>> <<<P000766>> формы
<\code{$n$<<\metavar{тип Аргументы}>. \id(\metavar{аргументы})}.
Если <<P000768> не обозначает класс <<P000769>
который объявляет конструктора по имени <<<P001519>,
Происходит ошибка времени компиляции.

<<<P001713>%
В противном случае, если <<P000771> происходит в постоянном контексте
(<<P001261>)
Затем <<<P000772> рассматривается как <<<P001520>,
Если <<P000774> Не происходит в постоянном контексте.
Затем <<<P000775> рассматривается как <<<P001521>.

% Мы можем добавить поддержку передачи аргументов типа статичным методам.
%, но это не та функция, которая появится в ближайшее время.
<<<P002017>{%>
Обратите внимание, что <<<P000777> Не может быть ничего, кроме создания экземпляра.
(постоянный или нет)
<<<P000778> приводит фактические аргументы типа <<<P000779>,
который не поддерживается, если <<P000780> обозначает префикс библиотеки,
и если <<P000781> Статический метод вызова.%
?
<<<P002018>

<<<P001714>%
<<<P002019><<<P001522> \code{<<P000783>>>-{}-}}
Рассмотрим выражение postfix <<P000784> формы <<P001523>,где <<P000786> является идентификатором и <<<P002021> является либо \lit{++} или \lit{-{}-}.
Ошибка времени компиляции возникает, если только <<P000787> обозначает переменную,
<<P000788> обозначает геттер и имеется связанный с ним сеттер <<<P001524>.
<<<P000790> Статический тип переменной <<<P000791> Возвратный тип геттера.
Ошибка времени компиляции возникает, если <<<P000792> не является <<<P002024>
<<P000793> У вас нет оператора <<<P002025>{+} (когда \op{} является \lit{++})
или оператор \lit{-} (когда \op{} является \lit{-{}-}),
или если тип возврата этого оператора не присваивается
переменная соответственно типу аргумента множителя.
Ошибка времени компиляции возникает, если <<<P001525> не может быть присвоено
Тип параметра указанного оператора.
Статический тип <<<P000794> является <<<P000795>.

<<<P001715>%
Оценка выражения postfix <<P000796>
Форма <<<P001526> соответственно \code{<<P000798>>>-{}-}
где <<<P000799> является идентификатором, действует следующим образом:
Оценить <<<P000800> на объект <<<P000801> Пусть <<<P000802> быть свежей переменной, связанной с <<<P000803>.
% % TODO(Eernst): Для того, чтобы исключить синтаксический сахар, мы должны
%% переписывают это, чтобы указать эффект непосредственно (случаи, чтобы помнить: 'v' - это
% % локальная переменная/параметр, экземпляр/статический/глобальный/импортный геттер.
Оценка <<<P001527> соответственно <<<P001528>.
<<<P000808> Показатель <<<P000809>.

<<<P002032>{%>
Вышесказанное гарантирует, что если оценка включает в себя геттер,
Он называется ровно один раз.
Аналогично в случаях ниже.%
?
<<<P002033>

<<<P001716>%
<<<P002034><<<P001529> <<<P002035><<<P000812>. <<<P000813>-{}-}}
Рассмотрим постфиксное выражение <<P000814> формы <<<P001530>,
где <<P000817> является буквальным и <<<P002036> является либо \lit{++}, либо <<P002038>>>{-{}-}.
Ошибка времени компиляции возникает, если только <<<P001531> обозначает статический геттер
и имеется связанный статический установщик <<<P001532>
<<<P002039>{(возможно, неявно вызванная статической переменной)}.
<<<P000821> Возвратный тип указанного геттера.
Ошибка времени компиляции возникает, если <<P000822> не является <<<P002040>
<<P000823> У вас нет оператора <<<P002041>{+} (когда \op{} является \lit{++})
или оператор <<<P002044>{-} (когда \op{} является \lit{-{}-}),
или если тип возврата этого оператора не присваивается
Тип аргументации сеттера.
Ошибка времени компиляции возникает, если <<<P001533> не может быть присвоено
Тип параметра указанного оператора.
Статический тип <<<P000824> является <<<P000825>.

<<<P001717>%
Оценка выражения postfix <<P000826>
Форма <<<P001534> соответственно <<<P002047><<<P000829>. <<P000830>-{}-}
где <<P000831> является типом буквальных поступлений следующим образом:
Оценить <<<P001535> на объект <<<P000834>
<<<P000835>> быть свежей переменной, связанной с <<P000836>.
Оценка <<<P001536> соответственно <<<P001537>.
<<<P000843> Показатель <<P000844>.
<<<P002048>

<<<P001718>%
<<<P002049>\code{<<<P000845>>>.<<<P000846>>>++} <<<P002050><<<P000847>. <<P000848>-{}-}}
Рассмотрим выражение postfix <<P000849> формы <<P001539>>>
где <<<P002051> <<<P002052>{++} или <<P002053>>>{-{}-}.
<<<P000852> Статический тип <<<P000853>.
Ошибка времени компиляции возникает, если <<<P000854> иметь
Геттер по имени <<P000855> и сеттер по имени <<<P001540>
<<<P002054>{(возможно, неявно вызванная переменной экземпляра)}.
<<<P000857> Возвратный тип указанного геттера.
Ошибка времени компиляции возникает, если <<P000858> не является <<<P002055>
<<P000859> У вас нет оператора <<<P002056>{+} (когда \op{} \lit{++})
или оператор <<<P002059>{-} (когда \op{} является \lit{-{}-}),
или если тип возврата этого оператора не присваиваетсяТип аргументации сеттера.
Ошибка времени компиляции возникает, если <<<P001541> не может быть присвоено
Тип параметра указанного оператора.
Статический тип <<<P000860> является <<<P000861>.

<<<<P001719>%
Оценка выражения postfix <<P000862>
Форма <<<P001542> соответственно <<<P002062><<<P000865>. <<P000866>-{}-}
идет следующим образом:
Оценить <<<P000867> на объект $u$ <<<P000869>> быть свежей переменной, связанной с <<<P000870>.
Оценить <<<P001543> на объект <<<P000873>
<<<P000874>> быть свежей переменной, связанной с <<<P000875>.
Оценка <<<P001544> соответственно <<<P001545>.
Далее <<P000882> Показатель <<P000883>.
<<<P002063>

<<<P001720>%
<<<P002064><<<P001546> \code{<<P000886>>>[<<P000887>>>-{}-}}
Рассмотрим выражение postfix <<P000888> формы <<P001547>>>
где <<<P002066> <<<P002067>{++} или <<P002068>>>{-{}-}.
<<<P000891> Статический тип <<<P000892>
<<P000893> Статический тип <<<P000894>.
Ошибка времени компиляции возникает, если <<<P000895> иметь
оператор \lit{[]} и оператор \lit{[]=}.
Пусть <<<P000896> будет возвратным типом первого.
Ошибка времени компиляции возникает, если только <<P000897> может быть присвоено
Первый тип параметра указанного оператора <<<P002071>{[]=}.
Ошибка времени компиляции возникает, если <<P000898> не является <<<P002072>
<<P000899> У вас нет оператора \lit{+} (когда \op{} является \lit{++})
или оператор <<<P002076>{-} (когда \op{} является <<P002078>>>{-{}-}),
или если тип возврата этого оператора не присваивается
Второй тип аргумента указанного оператора <<<P002079>{[]=}.
% Разрешается «e1[e2]++»; также, когда запись имеет двойной статический тип,
%, поэтому мы не можем просто сказать «назначаемый».
Ошибка времени компиляции возникает при прохождении целого числа <<<P001548>
как аргумент в пользу оператора \lit{+} или <<P002081>>>{-} Это было бы ошибкой.
Статический тип <<<P000900> является <<<P000901>.

<<<P001721>%
Оценка выражения postfix <<P000902>
Форма <<<P001549> соответственно <<<P002082><<<P000905> <<<P000906>-{}-}
идет следующим образом:
Оценить <<<P000907> на объект $u$ и <<<P000909> Объект <<<P000910>.
<<<P000911> и <<P000912> быть свежими переменными, связанными с <<<P000913> и <<<P000914> соответственно.
Оценить <<<P001550> на объект <<<P000917>
<<<P000918>> быть свежей переменной, связанной с <<<P000919>.
Оценка <<<P001551> соответственно <<<P001552>.
Далее <<P000926> Показатель <<P000927>.
<<<P002083>

<<<P001722>%
<<<P002084><<<P001553> <<<P002085><<<P000930>? <<P000931>-{}-}}
Рассмотрим выражение postfix <<P000932> формы <<P001554>
где <<<P002086> <<<P002087>{++} или <<P002088>>>{-{}-}.
Ошибки времени компиляции, которые могут быть вызваны
<<P001555>
Они также генерируются в случае <<<P001556>.
Статический тип <<<P000939> Статический тип <<<P001557>.

<<<P001723>%
Оценка выражения postfix <<P000942>
Форма <<<P001558> соответственно <<<P002089><<<P000945>? <<P000946>-{}-}
идет следующим образом:
Если <<P000947> Буквальный тип, оценка <<<P000948> является эквивалентным
Оценка <<<P001559> соответственно <<<P002090><<<P000951>. <<P000952>-{}-}.
В противном случае оценка <<P000953> к объекту <<P000954>.
Если <<P000955> является нулевым объектом, <<P000956> оценивает нулевой объект (<<<P001262>).
В противном случае пусть <<P000957> быть свежей переменной, связанной с <<<P000958>.
Оценка <<<P001560> соответственно <<<P002091><<<P000961>. <<P000962>-{}-} к объекту <<P000963>.
Далее <<P000964> Показатель <<P000965>.
<<<P002092>


<<<P002093> Соответствующие выражения
<<<P001779>

<<<P001724>%
<<<P002094> Назначаемые выражения - это терминыОн может появиться на левой стороне задания.
В этом разделе описывается, как оценивать термины этих терминов, когда это необходимо.
Семантика назначения в целом описывается в другом месте.
(<<P001263>).

<<<P002095>{%>
Грамматика присваиваемых выражений включает очень общие формы.
как идентификатор <<<P001602> Квалифицированный идентификатор <<<P001561>.
Следовательно, присваиваемое выражение может иметь множество различных значений,
В зависимости от связывания этих идентификаторов в контексте.

Например, термин <<<P001562> является присваиваемым выражением.
<<P001563> может относиться к геттеру <<<P001564> на
объект, на который ссылается переменная <<<P001565>,
В этом случае <<<P001566> будет оцениваться на предмет
Прежде чем перейти к <<<P001567> члена этого объекта.
Термин «P001568» может также обозначать класс <<<P001569> ссылаться через
Префикс импорта <<<P001570>,
<<<P001571> Статическая переменная этого класса.
В данном случае <<<P001572> не будет оцениваться в стоимостном выражении.%
?

<<P001194>
<assignableExpression> ::= <primary> <assignableSelectorPart
<<<P002096> \SUPER{} <безусловный Подписываемый Выборщик>
<<<P002098> <идентификатор>

<assignableSelectorPart> ::= <selector>* <assignableSelector>

<unconditionalAssignableSelector> ::= ‘[<выражение>] "
<<<P002099>> <идентификатор>

<assignableSelector> ::= <unconditionalAssignableSelector>
<<<<P002100>> <определитель>>
<<<P002101 >> '?' <выражение> '' "

<<<P001223>

<<<P001725>%
Раздел о назначениях
(<<P001264>)
определяет статический анализ и динамическую семантику
различные формы присвоения.
Каждый из этих случаев применим, когда указанные термины удовлетворяют
Заданные боковые условия
<<<P002102>{%>
(например, один случай формы <<<P001573>> <<<P000971>>
быть выражением, тогда как <<<P001574>
требует, чтобы <<<P000975> был буквальным типом)%
}
Случаи, требующие, чтобы термины были выражениями, считаются наименее конкретными.
Они используются только в том случае, если нет других совпадений.
<<<P002103>{%>
(содержащий <<P000976>) используется
если соответствующий термин является буквальным)%
}

<<<P002104>{%>
Синтаксически эти выражения не являются производными от <<<P002105>{выражение},
но они могут быть получены из <<<P002106>{выражение}, например, потому что
<<<<P002107>< Выражение может содержать <<<P002108>{первичное}
Выводится из <<<P002109>{выражение}.
Мы используем следующее правило, чтобы найти такие выражения:
?

<<<P001726>%
<<<P001787>%
Пусть <<<P000977> будет <<<P002110>{назначаемое выражение}.
Предположим, что <<P000978> Это такой термин, как <<P000979> могут быть получены из
\syntax{$t$ <assignableSelector>}.
В данном случае мы говорим, что <<P000981> Это
<<<<P002112>{приемный термин}{назначаемое выражение! термин "получатель"
<<P000982>.
Когда <<P000983> Это выражение, мы говорим, что <<P000984> Это
<<<P002113>{выражение приемника}{назначаемое выражение! выражение приемника
<<P000985>.

<<<P002114>{%>
Короче говоря, мы получаем <<<P000986> путем отсечения назначенного селектора
с конца <<P000987>.
Легко заметить, что только некоторые <<<P002115>
иметь срок приема.
Например, простой <<<P002116> нет.%
?

<<<P001727>%
Пусть <<<P000988> будет присваиваемым выражением.
Предположим, что <<P000989> имеет выражение приемника <<<P000990>.
Оценка <<<P000991> происходит таким же образом, как и
Оценка любого другого выражения.


<<<<P002117>{Лексический поиск}
<<P001780>

<<<P001728>%
В этом разделе указано, как искать имя
основанные на замкнутых лексических рамках.
Это называется P001298.
Когда <<<P001603> является идентификатором,
Вы можете найти имя <<P000992> Форма <<<P001604> Как и форма <<<P001575>.

<<<P001729>%
Лексический поиск может дать декларацию или префикс импорта.
И он ничего не может дать.<<<P002118>{%>
Это не ошибка времени компиляции, когда лексический поиск ничего не дает.
В этой ситуации данное имя <<<P000993> преобразуется в
<<<P001576>,
Статический анализ полученного выражения
может иметь или не иметь ошибок времени компиляции.%
?

<<<P001730>%
Лексический поиск может привести к ошибке времени компиляции.
как указано ниже.
Однако это отличается от результата, полученного в результате лексического поиска.
Потому что эта спецификация никогда не указывает на распространение ошибок.

<<<<P002119>{%>
Другими словами, когда другие части этой спецификации указывают, что
проводится лексический осмотр,
Они должны рассмотреть дальнейшие шаги, предпринятые
Когда всплывающее окно дает декларацию,
Когда он дает префикс импорта,
И когда он ничего не дает.
Но они не упоминают об этом, например, <<P001577> Это ошибка
потому что лексический поиск для <<<P001605> Погрешность.%
?

<<<P002120>{%>
Лексический поиск отличается от поиска члена экземпляра
(<<<P001265>)
потому что эта операция ищет через последовательность суперклассов,
В то время как лексический поиск ищет через последовательность замкнутых областей.
Лексический поиск отличается от простого поиска в закрытых областях
потому что лексический поиск «связывает» геттеров и сеттеров,
Подробнее ниже.%
?

<<<P001731>%
<<<P001788>%
Рассмотрим ситуацию, когда имя <<P000995>
Базовое имя <<<P001606>
(<<P001266>)
где <<<P001607> является идентификатором,
и лексический поиск <<<P000996> выполняется из определенного места
\DefineSymbol{\ell}.

<<<P002123>{%>
Мы указываем имя и место, откуда выполняется лексический поиск.
Месторасположение не всегда избыточное:
В некоторых ситуациях мы выполняем поиск сеттера по имени <<<P001578>,
Но токен <<P001579> В программе этого не происходит.
Чтобы справиться с такими ситуациями, мы должны указать оба.
имя, которое высматривается,
и местоположение, которое определяет, какие области являются закрытыми.%
?

<<<P001732>%
Когда мы говорим, что лексический поиск идентификатора <<<P001608> выполняется,
Понятно, что поиск осуществляется из
место данного происшествия <<<P002124>.

<<<P001733>%
<<<P000997> быть самой внутренней лексической областью, содержащей <<P000998>
который имеет декларацию с именем базы <<<P002125>.
В случае, когда <<P000999> иметь
Декларация под названием <<<P001609> а также заявление под названием <<<P001580>,
давать <<<P002126> D) быть декларацией под названием <<<P001000>.
В ситуации, когда <<<P001001> иметь
Одно заявление с именем basename <<<P002127>,
Пусть <<<P001002> будет такой декларацией.

<<<P002128>{%>
Нелокальная вариабельная декларация с именем <<<P001610> будет косвенно стимулировать
Геттер <<<P001611> и, возможно, сеттер <<<P001581> в текущем объеме.
Это означает, что <<<P001003> может обозначать неявно индуцированный геттер или сеттер
Вместо базовой переменной декларации.
Это важно в том случае, когда возникает ошибка.
потому что поиск был для одного вида, но только другой вид существует.
?

<<<P002129>{%>
Если мы ищем имя <<<P001004> с базовым именем <<<P002130>,
Мы прекращаем поиск, если находим какое-либо заявление под названием <<<P001612> или <<<P001582>.
Если в этой области имеются заявления как для <<<P001613>, так и для <<<P001583>,
Мы возвращаем тот, который имеет запрашиваемое имя <<<P001005>.
В случае, если присутствует только одно заявление, мы возвращаем его.
Даже если он может иметь имя <<<P001614>, когда <<<P001006> является <<<P001584>, или наоборот.
Эта ситуация может привести к ошибке, как указано ниже.
?

<<<P001734>%
На первом этапе мы проверяем несколько возможных ошибок.

% Даже если мы нашли заявление, оно может иметь неправильное название.
<<<P001735>%
<<<P002131><<<P001007> существует
В этом случае, по крайней мере, одно заявление с базовым именем <<<P001615>
находится в зоне действия по адресу <<<P001008>.Это ошибка времени компиляции, если имя <<<P001009> не является <<<P001010>,
Если только <<P001011> является членом экземпляра или локальной переменной
<<<P002132>{(который может быть формальным параметром)}.

<<<P002133>{%>
То есть, это ошибка, если мы ищем сеттера и находим геттера,
или наоборот,
Но не ошибка, если мы ищем сеттера и находим локальную переменную.
Если мы ищем сеттера и находим экземпляр геттера, или наоборот,
Это не ошибка,
Потому что сеттер может быть унаследован.
Это проверяется после того, как ничего не дает.
(что означает, что <<<P002134>{} будет подготовлено.%)
?

<<<P001736>%
Если <<<P001012> является членом инстанции,
Ошибка времени компиляции <<<P001013> Не имеет доступа к <<<P002135>.
<<P002136>

<<<<P001737>%
<<<P002137><<<P001014> не существует
Ошибка времени компиляции <<<P001015> не имеет доступа к <<<P002138>
(<<P001267>).
<<P002139>

<<<P002140>{%>
Мы всегда смотрим вверх <<<P002141> <<<P001616> и <<<P001585>,
Независимо от того, является ли <<<P001016> <<<P001617> или <<<P001586>.
Такой подход создает более тесную связь между парой деклараций.
где один является геттером по имени <<<P001618>
А другой — сеттер по имени <<<P001587>.
Это позволяет разработчикам думать о
геттер и сеттер, которые объявляются вместе как единое целое,
Вместо двух независимых деклараций.
?

<<<P002142>{%>
Например, если термин относится к <<<P001619> Нужен сеттер,
и самой внутренней декларации, названной <<<P001620> или <<<P001588> является геттером <<<P001017>
и нет соответствующего сеттера,
Это ошибка времени компиляции.
Эта ошибка возникает даже в том случае, когда более удаленная область охвата имеет
Декларация сеттера <<<P001018> <<<P001589>>
потому что мы уже взяли на себя обязательство использовать <<<P001019>
(так что это на самом деле «пара сеттера / геттера, где сеттер отсутствует»)
Можно сказать, что эта «парная» тень <<<P001020>:%
?

<<P000000>

<<<P001738>%
На втором и последнем шаге, если ошибки не произошло,
действовать, как описано в первом применимом случае из следующего перечня:

<<P001195>
<<P002143>
Когда <<<P001021> не существует,
Лексический поиск ничего не дает.
<<<P002144>{%>
В этом случае гарантируется, что <<<P001022> имеет доступ к <<<P002145>,
<<<P001023> будет рассматриваться как <<<P001590>.
Ошибки могут возникать, например,
потому что нет члена интерфейса <<<P002146> <<<P001025>,
а также нет доступного и применимого метода расширения.%
?
<<P002147>
Рассмотрим случай, когда <<<P001026> Формальное заявление о параметрах типа
класса или смеси.
Ошибка времени компиляции <<<P001027> происходит внутри
статический метод, статический геттер или статический сеттер,
внутри статического переменного инициализатора.
% NB: существует ошибка _no_, когда она возникает в инициализаторе переменной экземпляра;
%, что означает, что в данном конкретном случае разрешается «доступ к этому».
% Это может показаться непоследовательным, но это не должно быть вредным.
В противном случае лексический поиск дает <<P001028>.
<<P002148>
Рассмотрим случай, когда <<<P001029> является заявлением государства-члена в
Класс или смесь <<<P001030>.
Лексический поиск ничего не дает.
<<<P002149>{%>
В этом случае гарантируется, что <<<P001031> имеет доступ к <<<P002150>.%
?
<<P002151>
% Случаи, рассмотренные здесь: <<P001032> это
% - импорт с префиксом <<<P002152>;
% - псевдоним класса или типа; %, чье имя должно быть <<<P002153>, не нужно говорить, что.
% — библиотечная переменная, геттер или сеттер;
% - статический метод/настройщик/настройщик класса, смеси, перечня или расширения;
% - примерный способ расширения;
% — локальная переменная (которая может быть формальным параметром);
% — локальная функция.
В противном случае лексический поиск дает <<<P001033>.
<<<P001224>

<<<P002154>{%>Обратите внимание, что лексический поиск никогда не даст объявление о том, что
Член инстанции.
В каждом случае, когда установлено, что ошибки нет,
и один из членов экземпляра является результатом поиска,
Лексический поиск ничего не дает.

Причина этого заключается в том, что не может быть декларации в объеме,
но интерфейс класса может иметь требуемый член,
Возможно применение метода расширения.
Для единообразия,
В таких ситуациях мы всегда сообщаем, что ничего не нашли.
Это означает, что <<<P002155> Должен быть добавлен.%
?


<<<P002156> Ссылка на идентификатор
<<P001781>

<<<P001739>%
<<<P001299> состоит из одного идентификатора;
Он обеспечивает доступ к объекту через неквалифицированное имя.
<<<P002157>{typeIdentifier} - идентификатор, который может быть использован
как название вида декларации.

<<<P002158>{%>
A <<<P002159>{qualifiedName} не является выражением идентификатора,
Но мы определяем его синтаксис здесь, потому что он используется в нескольких различных контекстах.
и более тесно связана с простым идентификатором
чем для любого из тех правил грамматики, где он используется.
?

<<P001196>
<identifier> ::= <Идентификатор>
<<P002160> <BUILT<<<P002352> В <<<P002353> Идентификатор
<<P002161> <Другой<<<P002354> Идентификатор

<typeIdentifierNotType> ::= <IDENTIFIER>
<<P002162> <Другой<<<P002355> Идентификатор <<<P002356> <<<P002357> Тип
<<P002163> \DYNAMIC{}

<typeIdentifier> ::= <typeIdentifierNotType>
<<P002165> \TYPE{}

<qualifiedName> ::= <typeIdentifier> <identifier>
<<<P002167> <typeIdentifier> <typeIdentifier> <identifier>

<BUILT<<<P002358> В <<<P002359> Идентификатор: := \ABSTRACT{} | <<P002169>>>{} | <<P002170>>>{} | <<P002171>>>{}
\alt\hspace{-3mm}<<P002174>>>{} |<<P002175>>>{} |<<P002176>>>{} |<<P002177>>>{} !
\FACTORY{} |<<P002179>>>{} |<<P002180>>>{}
\alt\hspace{-3mm}<<P002183>>>{} |<<P002184>>>{} |<<P002185>>>{} |<<P002186>>>{} !
\LIBRARY{} |<<P002188>>>{} |<<P002189>>>{}
\alt\hspace{-3mm} \PART{} | \REQUIRED{} | \SET{} | \STATIC{} | \TYPEDEF{}

<Другой<<<P002360> Идентификатор <<<P002361> Не <<<P002362> Тип ::= \gnewline{}
\ASYNC{} | <<P002199>>>{} | <<P002200>>>{} | <<P002201>>>{} | <<P002202>>>{} | <<P002203>>>{} !
\AWAIT{} | \YIELD{}

<Другой<<<P002363> Идентификатор: := \gnewline{}
<Другой<<<P002364> Идентификатор <<<P002365> <<<P002366> TYPE> | \TYPE{}

<IDENTFIER<<<P002367> NO <<<P002368> ДОЛЛАР: := <IDENTFIER<<<P002369> Начало<<<P002370> No<<<P002371> ДОЛЛАР
\gnewline{} <IDENTFIER<<<P002372> Часть <<<P002373> No<<<P002374> Доллар*

<IDENTFIER<<<P002375> Начало <<<P002376> No<<<P002377> DOLLAR > ::= <LETTER> "

<IDENTFIER<<<P002378> Часть <<<P002379> No<<<P002380> ДОЛЛАР: := \gnewline{}
<IDENTFIER<<<P002381> Начало <<<P002382> No<<<P002383> ДОЛЛАР > | <ПРИЛОЖЕНИЕ

<IDENTIFIER> ::= <IDENTFIER<<<P002384> Начало> <IDENTIFIER<<<P002385> Часть*

<IDENTFIER<<<P002386> Начало>::= <IDENTFIER<<<P002387> Начало <<<P002388> NO<<<P002389> ДОЛЛАР "

<IDENTFIER<<<P002390> Часть ::= <IDENTFIER<<<P002391> Начало>> <Цифры>

<LETTER> ::= 'a' .. 'z' | 'A' .. 'Z'

<DIGIT> ::= 0.. "9"

<WHITESPACE> ::= (<<<P002392><<P002210>> | <LINE<<<P002393>> BREAK > +
<<<P001225>

<<<P001740>%
Упорядочение лексических правил выше гарантирует, что <<<P002211>
<<<P002212> Идентификатор <<<P002394> NO<<<P002395> DOLLAR не выводят никаких встроенных идентификаторов.
Аналогично, лексическое правило для зарезервированных слов (\ref{reservedWords})
должно считаться предшествующим правилу для <<<P002213> BUILT<<<P002396> В <<<P002397> Идентификатор.
такой, что \synt{IDENTIFIER}} и \synt{ Идентификатор <<<P002398> No<<<P002399> ДОЛЛАР
Также не следует выводить какие-либо зарезервированные слова.

<<<P001741>%
<<<P001300> является одним изидентификаторы, полученные в результате производства <<<P002216> BUILT<<<P002400> В <<<P002401> Идентификатор.

<<<P002217>{%>
Обратите внимание, что это синтаксическая ошибка, если встроенный идентификатор
используется как заявленное название
% % TODO(Eernst): Приходите с расширениями, добавьте их.
префикс, класс, микширование, число, параметр типа, псевдоним типа или расширение.
Аналогично, это синтаксическая ошибка для использования встроенного идентификатора.
кроме других \DYNAMIC{} или <<P002219>>>{}
как идентификатор в аннотации типа или связанный параметр типа.%
?

<<<P002220>{%>
Встроенные идентификаторы — это идентификаторы, которые используются в качестве ключевых слов в Dart.
Но это не зарезервированные слова.
Встроенный идентификатор не может использоваться для обозначения класса или типа.
Другими словами, они рассматриваются как зарезервированные слова при использовании в качестве типов.
Это устраняет многие запутанные ситуации.
как для читателей, так и во время парсинга.%
?

<<<P001742>%
Ошибка времени компиляции, если один из идентификаторов \AWAIT{} или <<P002222>>>{}
используется в качестве <<<P002223>{идентификатора} в функциональном теле
отмеченный либо <<<P002224>, <<P001591> или <<P001592>.

<<<<<<<<<<<<<<<<<<<<<<<<<<<<<<<<<<<<<<<<<<<<<<<<<<<<<<<<<<<<<<<<<<<<<<<<<<<<<<<<<<<<<<<<<<<<<<<<<<<<<<<<<<<<<<<<<<<<<<<<<<<<<<<<<<<<<<<<<<<<<<<<<<<<<<<<<<<<<<<<<<<<<<<<<<<<<<<<<<<<<<<<<<<<<<<<<<<<<<<<<<<<<<<<<<<<<<<<<<<<<<<<<<<<<<<<<<<<<<<<<<<<<<<<<<<<<<<<
Это делает идентификаторы \AWAIT{} и \YIELD{} вести себя как зарезервированные слова
в ограниченном контексте.
Этот подход был выбран потому, что он был менее разрушительным, чем это было бы.
делать \AWAIT{} и \YIELD{} зарезервированные слова или встроенные идентификаторы;
В то время, когда эти функции были добавлены к языку.
?

<<<P001743>%
<<<P001301> Это два или три идентификатора, разделенные <<<P002230>{.}.
Все, кроме последнего, должно быть <<<P002231>{typeIdentifier}.
Он используется для обозначения декларации, которая импортируется с префиксом.
или заявление <<<P002232>{} в виде класса, смеси, перечня или расширения, или обоих.

<<<P001744>%
Статический тип выражения идентификатора <<<P001034> который является идентификатором <<<P001621>
определяется следующим образом.
Выполните лексический поиск <<<P001622>
(перенаправлено с «P001269»)
Местонахождение <<<P001035>.

<<<P001745>%
<<<P002233> Лексический поиск дает декларацию.
<<<P001036> быть декларацией, полученной в результате лексического поиска <<<P002234>.

<<P001197>
<<P002235>
Если <<<P001037> объявляет класс, миксинг, энум, псевдоним типа, перечисленный тип,
или параметра типа,
Статический тип <<<P001038> является <<<P001593>.
<<P002236>
Если <<P001039> Декларация библиотечного геттера
<<<P002237>{(которые могут быть неявно вызваны библиотечной переменной)},
Статический тип <<<P001040> является статичным типом
библиотечный геттер вызов <<<P001623>
(<<P001270>).
<<P002238>
Если <<P001041> Статический метод, библиотечная функция или локальная функция.
Статический тип <<<P001042> Тип функции <<<P001043>.

<<<P002239>{%>
Обратите внимание, что <<<P001044> может быть впоследствии подвергнут
обобщение функции
(<<<P001271>)%
?
<<P002240>
Если <<<P001045> Объявление статического геттера
<<<P002241>{(которые могут быть индуцированы статической переменной)}
<<<P001046> встречается в классе <<<P001047>,
Статический тип <<<P001048> Тип возврата Геттера
<<<P001594>.
<<P002242>
Если <<<P001050> локальная переменная декларация
<<<P002243>{(который может быть формальным параметром)}
Статический тип <<<P001051> тип переменной <<<P001052> Заявлено <<<P001053>,
Если только <<<P001054> Известно, что имеет некоторый тип <<<P001055>,
где <<<P001056> является подтипом любого другого типа <<<P001057>
Таким образом, <<<P001058> Известно, что имеет тип <<<P001059>,
В этом случае статический тип <<<P001060> является <<<P001061>.
<<P002244>
% Продление геттера вызова по объему.
Если <<<P001062> Декларация экземпляра Getter
в заявлении о продлении <<<P001063>,
Статический тип <<<P001064> Тип возврата <<<P001065>.
<<P002245>
% Отрыв метода расширения по объему.Если <<<P001066> является декларацией метода экземпляра
в декларации о продлении <<<P001067>,
Статический тип <<<P001068> Тип функции <<<P001069>.
<<P002246>
<<<P002247>{%>
Лексический взгляд никогда не даст декларации.
который является членом класса.%
?
<<<P001226>
\vspace{-1ex)
<<P002249>

<<<P001746>%
<<<P002250>> Лексический поиск дает префикс импорта
В этом случае лексический поиск
(<<<P001272>)
<<<P001624> Дает префикс импорта <<P001070>.
Префикс никогда не может быть использован в качестве отдельного выражения.
В этом случае возникает ошибка времени компиляции,
Если токен сразу после <<<P001071> <<<P002251>{.}.
Статический тип не связан с <<<P001072> в данном случае.

% Тип указывается в <<<P001273>, поэтому было бы последовательно сообщать
%. Однако из этого типа не будет использоваться, поскольку статический тип
% от всех построений формы <<<P001073>.Что-то указано без ссылки
% к статическому типу <<<P001074>. Поэтому мы не упоминаем этот тип здесь.
<<<P002252>{%>
Такой тип не нужен, потому что каждая конструкция
Префикс импорта <<<P001075> используется и сопровождается <<<P002253>{.} указывается
таким образом, что тип <<<P001076> Не используется.%
?
<<P002254>

<<<P001747>%
<<<P002255> Лексический поиск ничего не дает
Когда лексический взгляд
(<<<P001274>)
<<<P001625> ничего не дает,
<<P001077> рассматривается как
(<<<P001275>)
<<P001595>.

<<<P002256>{%>
В данном случае известно, что <<<P001078> Имеет доступ к <<<P002257>
(<<<P001276>).
Статический анализ и оценка продолжаются
<<P001596>,
поэтому нет необходимости дополнительно указывать лечение <<<P001079>.%
?
<<P002258>

<<<P001748>%
Оценка выражения идентификатора <<<P001080> формы <<<P001626>
идет следующим образом:

<<<P001749>%
<<<P002259> Лексический поиск дает декларацию.
В этом случае лексический поиск
(перенаправлено с «P001277»)
<<<P001627> выдает декларацию <<<P001081>.
Оценка <<<P001082> идет следующим образом:

<<P001198>
<<P002260>
Если <<<P001083> является классом, микшированием, enum или типом alias,
Значение <<<P001084> является объектом, реализующим класс <<<P001597>
который соответствует соответствующему типу.
<<P002261>
Если <<<P001085> Типовой параметр <<<P001086> Тогда значение <<<P001087> это
значение фактического типа аргумента, соответствующего <<<P001088>
Это было передано созидателю, который создал
Текущее связывание <<<P002262>.
<<P002263>
Если <<<P001089> Декларация библиотечного геттера
<<<P002264>{(которые могут быть неявно вызваны библиотечной переменной)},
Оценка <<<P001090> равносильно оценке призыва к
Библиотечный геттер <<<P001628>
(<<P001278>).
<<P002265>
% Мы могли бы сказать, что «выталкивание вызвано постоянной переменной»; но всегда будет
% - переменная, потому что декларация Геттера не может дать константу, поэтому мы
% используйте ярлык и не упоминайте геттер полностью здесь.
Если <<P001091> библиотека, класс или локальная постоянная переменная одной из форм
\code{<<P002267>>>{) $v$ = $e'$;} или \code{\CONST{} <<<P001094> <<<P001095> = <<<P001096>;
Тогда значение <<<P001097> значение постоянного выражения <<<P001098>.
<<P002270>
Если <<P001099> является декларацией
функция верхнего уровня, статический метод или локальная функция;
<<<P001100> оценивает функцию объекта, полученную путем клозулизации
(<<<P001279>)
<<<P001101>.
<<P002271>
Если <<<P001102> является заявлением на получение экземпляра в заявлении о продлении <<<P001103>,
<<<P001104> оценивает результат вызова указанного геттера
с текущим связыванием <<<P002272>, и
текущие связывания параметров типа, объявленные по <<<P001105>.
<<P002273>Если <<p001106> является заявлением метода экземпляра в заявлении о расширении <<<P001107>
с параметрами типа <<<P002274>{X}{1}{}
<<<p001108> оценивает результат клозуризации метода расширения
\code{<<P001109>>>\List{X}{1}{s}>(<<P002277>>>). <<<P002278>>
(<<<P001280>).
<<P002279>
Если <<p001110> локальная переменная <<<P001111>
<<<P002280>{(который может быть формальным параметром)}
<<<P001112> оценивает текущее связывание <<<P001113>.
<<P001227>

<<<P002281>{%>
Обратите внимание, что <<<P001114> не может быть декларацией
статическая переменная, статический геттер или статический сеттер, объявленный в классе <<<P001115>,
Потому что в этом случае <<<P001116> рассматривается как
(<<P001281>)
Добыча имущества
(<<P001282>)
<<<P001598>,
который также определяет оценку <<<P001118>.%
?
<<P002282>

<<<P001750>%
<<P002283> Лексический поиск дает префикс импорта
Такая ситуация не может возникнуть,
Потому что это ошибка времени компиляции
для оценки префикса импорта как выражения,
и никаких конструкций с импортным префиксом
\commentary{(например, извлечение свойств <<P001599>>>)}
Оцените префикс импорта.
<<P002285>

<<<P001751>%
<<<P002286> Лексический поиск ничего не дает
Такая ситуация не может возникнуть,
Это происходит только тогда, когда <<<P001121> рассматривается как
<<<P001600>,
чья оценка указана в другом месте
(<<P001283>).
<<P002287>


<<<P002288> Тип испытания
<<P001782>

<<<P001752>%
<<<P001302> Тестирование, если объект является членом типа.

<<P001199>
<typeTest> ::= <isOperator> <typeNotVoid>

<isperator> ::= <<<P002289>{}!
<<<P001228>

<<<P001753>%
Оценка выражения <<<P002290><<<P001122> <<<P002291>< <<<P001123>> идет следующим образом:

<<<P001754>%
Выражение <<<P001124> оценивается на объект <<<P001125>.
Если динамический тип <<<P001126> является подтипом <<<P001127>,
Показатель выражается в <<<P002292>.
В противном случае он оценивает до <<<P002293>.

<<<P002294>{%>
Из этого следует, что <<<P002295>{$e$ \IS{} Объект всегда верен.
Это имеет смысл в языке, где все является объектом.

Также обратите внимание, что \code{\NULL{} \IS{} <<<P001129>> ложный
Если только <<P001130> <<<P001131> или <<P001132>.
Первые два бесполезны, как и все остальное.
форма \code{$e$ \IS{} Объект или <<<P002302><<<P001134> \IS{} <<<P002304>>.
Пользователи должны проверить нулевой объект (<<P001284>) непосредственно
а не с помощью типовых тестов.%
?

<<<P001755>%
Выражение <<<P002305><<<P001135> <<<P002306>{}! <<<P001136>> является эквивалентным
<<<P002307><! (<<P001137>) \IS{} <<<P001138>>.

<<<P001756>%
<<<P001139> Быть локальной переменной
<<<P002309>{(который может быть формальным параметром)}.
Выражение формы \code{$v$ \IS{} <<<P001141>>
Показывает, что <<<P001142> имеет тип <<<P001143>
Если <<P001144> является подтипом типа выражения <<<P001145>.
%
В противном случае
если заявленный тип <<<P001146> переменная типа <<<P001147>,
<<<P001148> является подтипом границы <<<P001149>,
<<<P001150> является подтипом типа выражения <<<P001151>,
<<<P001152> Показывает, что <<<P001153> Имеет тип <<<P001154>.
%
В противном случае <<P001155> Не показывает, что <<<P001156> имеет тип <<<P001157> для любого <<<P001158>.

<<<P002312>{%>
Мотивацией для «показов, что v имеет отношение типа T» является
Чтобы уменьшить ложные ошибки, тем самым обеспечивая более естественный стиль кодирования.
Правила в текущей спецификации намеренно держатся простыми.
В будущем было бы целесообразно усовершенствовать эти правила;
Такое усовершенствование позволило бы принимать больше кода без ошибок.
Но не отказывайтесь от кода без ошибок.

Правило касается только местных жителей и параметров.Нелокальные переменные могут быть изменены
с помощью функций или методов, которые не доступны
для локального анализа.

Бессмысленно выводить более слабый тип, чем тот, который уже известен.
Кроме того, это приведет к ситуации, когда несколько типов
с переменной в заданной точке,
Это усложняет спецификацию.
Отсюда требование, что продвигаемый тип является подтипом текущего типа.

В любом случае, это не ошибка, когда тип теста не показывает
что данная переменная не имеет «лучшего» типа, чем ранее известный,
Но инструменты могут дать подсказку в таких случаях.
если подходящие эвристики указывают, что продвижение, вероятно, будет предназначено.%
?

<<<P001757>%
Статический тип является выражением <<<P001601>.


<<<P002313>> Тип Каста
<<P001783>

<<<P001758>%
<<<P001303> Это гарантирует, что объект является членом типа.

<<P001200>
<typeCast> ::= <asOperator> <typeNotVoid>

<asOperator> ::= \AS{}
<<P001229>

<<<P001759>%
Оценка литого выражения \code{$e$ \AS{} <<<P001160>> идет следующим образом:

<<<P001760>%
Выражение <<<P001161> Оценивается на объект <<<P001162>.
% Эта ошибка может произойти по замыслу «как».
Ошибка динамического типа, если <<P001163> не является нулевым объектом (<<<P001285>),
и динамический тип <<<P001164> не является подтипом <<<P001165>.
В противном случае <<P001166> Показатель <<<P001167>.

<<<P001761>%
Статический тип литого выражения \code{$e$ \AS{} <<<P001169>> <<<P001170>.


<<<P002319> Заявления
<<P001784>

<<<P001762>%
<<<P001304> Это фрагмент кода Dart, который может быть выполнен во время выполнения.
Заявления, в отличие от выражений, не оцениваются на предмет,
Но вместо этого они выполняются для их влияния на состояние программы и поток управления.

<<<P001201>
<statements> ::= <statement> *

<statement> ::= <label>* <nonLabelled Statement>

<nonLabelled Statement> ::= <блокировать>
<<<P002320> <localVariableDeclaration>
<<<P002321> <for Statement>
<<<P002322> <в то время как заявление>
<<<P002323> <doStatement>
<<<P002324 >> <переключатель
<<<P002325> <если заявление>
<<<P002326> <rethrow Statement>
<<<P002327> <try Statement>
<<<P002328> <break Statement>
<<P002329> <продолжение заявления>
<<<P002330> <возврат заявления>
<<P002331> <yieldStatement>
<<P002332> <yieldEachStatement>
<<<P002333> <выражение>
<<<P002334> <assert Statement>
<<<P002335> <localFunctionDeclaration>
<<<P001230>


<<<P002336> Заявление о завершении
<<P001785>

<<<P001763>%
Исполнение заявления \IndexCustom{завершение}{завершение}
одним из пяти способов:
либо
<<<P002338>{завершается нормально}{завершается нормально},
<<<<P002339>{breaks}{завершение!breaks}
Или это <\IndexCustom{продолжение}{завершение!продолжение}
(либо на этикетке, либо без этикетки),
<<<P002341>{возвращается}{завершение!возвращается} (с объектом или без него),
Или это <<<P002342>{броски}{завершение!броски}
Исключительный объект и связанный с ним след стека.

<<<P001764>%
В описаниях выполнения заявления по умолчанию является то, что
исполнение обычно завершается, если не указано иное.

<<<P001765>%
Если исполнение заявления, <<<P001171>,
определяется с точки зрения выполнения другого заявления,
и исполнение этого другого заявления обычно не завершается,
если не указано иное, исполнение <<<P001172> остановки
В этот момент и завершается одинаково.
<<<P002343>{%>
Например, если исполнение тела a <<<P002344> петля возвращает объект,
Так же выполняется выполнение <<<P002345>{} петлевое высказывание.%
?

